\section{Resolution map for the reported difficulties}
\label{sec:resolution-map}

The repository audit revealed a stable list of recurring obstacles.  The table below
records what the framework resolves, what it resolves only under explicit data, and
what remains substantive mathematics.  ``Resolved'' means that the generic reduction
is proved; it does not mean that every application packet has been migrated to a
kernel-checked object proof.

\small
\begin{longtable}{@{}p{0.16\textwidth}p{0.23\textwidth}p{0.32\textwidth}p{0.20\textwidth}@{}}
\toprule
Obstacle & Invalid or incomplete move & FRNF repair & Status / remaining input \\
\midrule
\endfirsthead
\toprule
Obstacle & Invalid or incomplete move & FRNF repair & Status / remaining input \\
\midrule
\endhead

Unbounded precision induction & Put a predicate quantifying over all rows of a Cauchy
name inside induction. & Freeze the final demand, compute the full dual-horizon closure,
materialize its finite table, and Booleanize all table contracts. & Generic reduction
resolved; object table library still to implement. \\
\addlinespace

Adaptive oracle use & Let future queries depend on higher-type answers without a fixed
finite domain. & Defunctionalize to a bounded transition system; compute worst-case
read and proof horizons; verify the actual trace. & Resolved when an explicit clock and
query majorant are supplied. \\
\addlinespace

Computation versus proof horizon & Track only rows read by the executable routine. &
Carry separate read and justification functions and propagate both backward. & Resolved
architecturally and mathematically. \\
\addlinespace

Exact boundary or chart choice & Decide sign, equality, or chart membership for an
arbitrary represented real. & Use overlapping branches or finite candidate clouds;
collapse only from a positive gap. & Conditional resolution; gap required for exact
label. \\
\addlinespace

Nonzero to inverse & Treat \(x\ne0\) as a rational lower bound. & Guard inverse by
\(\abs{x}\ge\lambda>0\), or clear denominators and prove a polynomial identity. &
Requires apartness input unless reformulated. \\
\addlinespace

Abstract compactness & Invoke compactness, attainment, or a finite subcover without
quantitative data. & Certified finite nets, coverage witnesses, and moduli; return
approximate extrema. & Conditional resolution; net/modulus producer required. \\
\addlinespace

Root and eigenvector labels & Globally order branches and vectors through
multiple roots or eigenspaces. & Return clusters, symmetric observables, projective
objects, or isolated spectral subspaces. & Resolved for invariant outputs; individual
label needs separation. \\
\addlinespace

Metastability to convergence rate & Infer an eventual tail modulus from one stable
finite window. & Keep the metastable contract; derive a Cauchy name only from a tail or
summable-motion modulus. & Exact boundary retained. \\
\addlinespace

Cardinality to cutoff & Stop after a number of found solutions equal to the claimed
cardinality. & Require a height/cutoff producer or a complete-list certificate. & No
generic operator; impossible by \cref{barrier:cardinality-cutoff}. \\
\addlinespace

Host Lean theorem to iPRA & Call a theorem about \(\N\), \(\Q\), or \(\mathbb R\) an
object proof. & Formula firewall: concrete PR syntax, object formula, local
\(\Prf\,\PRAAx\) constructor, proof schedule. & Engineering target; logical kernel
combinators supplied. \\
\addlinespace

Universal PR evaluator & Execute an arbitrary primitive-recursive program code inside
one PR function. & Statically specialize fixed evaluators, use finite tables, or a
restricted fuel-bounded interpreter. & Exact impossibility by diagonalization. \\
\addlinespace

Arbitrary Cauchy Puiseux data & Return exact leading exponent/coefficient from an
approximate coefficient name. & Require finite exact algebraic splitting/rank data, or
return residual clusters/regular branches. & Exact-output interface obstructed by the
zero-test barrier. \\
\addlinespace

Qualitative differentiability or holomorphy & Extract a uniform derivative or analytic
modulus from a bare qualitative assertion. & Supply a modulus or derive one from the
special presentation; then use finite partitions/series and explicit remainders. &
Conditional resolution. \\
\addlinespace

General Siegel/Faltings & Treat finite verification as a source of a height bound. &
Compile ``height bound implies complete list''; keep the bound producer visible. &
Finite consequence resolved; general producer remains deep mathematics. \\
\addlinespace

Strong PNT/uniform prime distribution & Expect demand bookkeeping to improve a weak
zero-free region or remove exceptional zeros. & Express every analytic estimate as a
boundary contract and preserve parameter dependence. & Formalization improved; stronger
analytic theorem remains. \\
\addlinespace

General probability laws & Quantify over arbitrary measures/functions and use weak
compactness. & Finite rational laws plus explicit quantization, tail, Lindeberg, and
smoothing moduli. & Conditional resolution; domain estimates required. \\
\addlinespace

Quotient equality & Rely on opaque host quotients and extensional equality. & Canonical
normal forms or certified setoids with finite equality witnesses. & Resolved as an
interface discipline. \\
\addlinespace

Infinite branch selection & Apply K\"onig's lemma to compatible finite levels. & Finite
backward viability for each horizon; coherence from shrinking diameter and motion
budgets. & Resolved for quantitative represented outputs. \\
\bottomrule
\end{longtable}
\normalsize

\subsection{Resolved does not mean automatic}

For example, the dual-horizon theorem gives a correct composition rule, but an
application still has to derive its local horizons.  The finite-net theorem gives an
approximate minimum once a net and modulus exist, but it does not construct those for
an arbitrary compact object.  The proof schedule builds a final derivation once local
\(\Prf\,\PRAAx\) proofs are supplied, but it does not synthesize domain lemmas from a
host proof.  The framework turns vague gaps into finite, typed obligations.

\subsection{Highest-value migration targets}

The audit suggests the following priority order.

\begin{enumerate}[label=\arabic*.,leftmargin=2.2em]
\item Complete the generic bounded-formula and encoded-table object library.
\item Formalize reverse demand closure and the two-pass viability algorithm as actual
\(\Prf\,\PRAAx\) theorems.
\item Migrate one representative analytic packet whose moduli are already explicit,
so that dual horizons are exercised end to end.
\item Migrate one relational application with collisions, such as stable root clusters
or simplicial carriers.
\item Express a conditional height-to-list theorem and connect a special-family height
producer.
\item Only then generalize the code generator to larger packet families.
\end{enumerate}

These tasks maximize reuse and directly attack the semantic-Lean-to-object-iPRA gap
identified throughout the corpus.
